% Part: sets-functions-relations
% Chapter: induction
% Section: induction

\documentclass[../../include/open-logic-section]{subfiles}

\begin{document}

\olfileid{sfr}{ind}{int}

\olsection{परिचय}

\begin{explain}
विगमन ही तर्कशास्त्रातील गणितात वारंवार वापरली जाणारी
सिद्धतेची पद्धत आहे. नैसर्गिक संख्यांवर केलेले विगमन
तुम्हाला गणितातून परिचित असेल. प्रत्यक्षात अनेक प्रकारचे
गणिती व तार्किक वस्तुसंच विगमनाला अनुकूल असतात.

सर्वसाधारणपणे, विगमनाने सार्वत्रिक दावा सिद्ध करता येतो;
म्हणजे एखाद्या प्रकारच्या प्रत्येक वस्तूला एक गुणधर्म आहे हे
दाखवता येते. विशेषतः, “सार्वत्रिक प्रस्तावना” या सिद्धता-पद्धतीने
सिद्धता करणे फार गुंतागुंतीचे असेल तेव्हा विगमन उपयुक्त ठरते.
विशिष्ट रीतीने बनवलेल्या गणिती वस्तूंवरच विगमन चालते:
संचातील प्रत्येक घटक मूलभूत असेल किंवा मूलभूत घटकांपासून
बनलेला असेल, तर त्या संचावर विगमन करता येते. अधिक नेमकेपणाने:
\end{explain}

\begin{defn}
एखादा संच~$S$ हा फलन~$f$ च्या खाली \emph{बंद} असतो तेव्हाच
आणि तेव्हाच $x \in
S$ असताना $f(x) \in S$ असते.
\end{defn}

\begin{explain}
उदाहरणार्थ, नैसर्गिक संख्यांचा संच ($\mathbb{N}$) उत्तरवर्ती
फलन $f(x) = x+1$ च्या खाली बंद आहे: $x$ नैसर्गिक संख्या
असेल, तर $x+1$ सुद्धा नैसर्गिक संख्या असते. मात्र नैसर्गिक
संख्यांचा संच पूर्ववर्ती फलन $f(x) = x-1$ च्या खाली बंद
नाही; कारण $0$ नैसर्गिक संख्या आहे, पण $-1$ नाही.
\end{explain}

\begin{defn}
फलन~$f(x)$ च्या प्रांतातील कोणत्याही वस्तू~$a$ साठी $P(a)$
असल्यास $P(f(a))$ असेल (म्हणजे $a$ ला गुणधर्म~$P$ असेल,
तर $f(a)$ लाही असेल), तर गुणधर्म~$P$ फलन~$f$ च्या खाली
\emph{जपला जातो} असे म्हणतात.
\end{defn}

\begin{explain}
“सम संख्या असणे” हा गुणधर्म $f(x) = x+2$ या फलनाखाली
जपला जातो; कारण $x$ सम असेल, तर $x+2$ सुद्धा सम असतो.
मात्र हा गुणधर्म उत्तरवर्ती फलनाखाली जपला जात नाही;
कारण $x$ सम असेल, तर $x+1$ विषम असतो.

नैसर्गिक संख्यांवर विगमन चालते, कारण प्रत्येक नैसर्गिक
संख्येपर्यंत 0 पासून उत्तरवर्ती फलन सान्त वेळा लावून
पोहोचता येते. ज्या संचावर विगमन करता येते, त्या प्रत्येक
संचाचे हे वैशिष्ट्य असते.
\end{explain}

\begin{defn}[$\Nat$ वरील विगमन]
0 ला गुणधर्म~$P$ असेल आणि $P$ हा गुणधर्म उत्तरवर्ती
फलनाखाली जपला जात असेल, तर प्रत्येक नैसर्गिक संख्येला
गुणधर्म~$P$ असतो.
\end{defn}

\begin{ex} 0 पासून $n$ पर्यंतच्या नैसर्गिक
संख्यांची बेरीज $\frac{n(n+1)}{2}$ असते. \\

\textbf{विगमनाचा आधार.} पहिल्या 0 नैसर्गिक
संख्यांची बेरीज $0=
\frac{0\cdot 1}{2}$ असते.\\

\textbf{विगमनाची पायरी.} गुणधर्म $k$ साठी खरा आहे असे
माना. तो $k+1$ साठीही खरा आहे हे दाखवू. म्हणजे $P(k)$
खरे आहे असे माना; अर्थात,
\[ 0 + 1 + \ldots + k = \frac{k(k+1)}{2} \]
$P(k+1)$ खरे आहे, म्हणजे
\[ 0 + 1 + \ldots + k + (k+1) = \frac{(k+1)(k+2)}{2} \]
हे $P(k)$ या गृहीतकाचा उपयोग करून दाखवू.

\begin{align*}
0 + 1 + \ldots + k &= \frac{k(k+1)}{2} \tag{गृहीतक} \\
0 + 1 + \ldots + k + (k+1) &= \frac{k(k+1)}{2} + (k+1)
\tag{दोन्ही बाजूंना $k+1$ मिळवा} \\
&= \frac{k(k+1) + 2(k+1)}{2} \\
&= \frac{k^2 + k + 2k +2}{2} \\
&=\frac{(k+1)(k+2)}{2}
\end{align*}
अशा रीतीने विगमनाची पायरी सिद्ध होते आणि दावा सिद्ध झाला.
\end{ex}

\begin{explain}
!!{formula}s बाबतीत मूलभूत घटक अण्वीय !!{formula}s असतात.
कमी गुंतागुंतीची !!{formula}s संकारकांनी जोडून अधिक गुंतागुंतीची
!!{formula}s मिळतात. प्रत्येक संकारकाला !!{formula}s वरचे
असे एक फलन अनुरूप मानता येते की ते संबंधित संकारक वापरून
दोन सूत्रे जोडणारे नवे !!{formula} देते. उदाहरणार्थ,
$!A$ आणि $!B$ ही दोन !!{formula}s जोडून संयुती
बनवणारे फलन $\mathcal E _\land (!A,
!B) = !A \land !B$ असे लिहिता येते. यांना
\emph{सूत्ररचना-फलने} म्हणू. !!{formula}s चा संच या
फलनांच्या खाली बंद आहे: $!A$ आणि $!B$ ही
!!{formula}s असतील, तर $\lnot !A$,
$!A \land
!B$ इत्यादीही तशीच असतात.
\end{explain}

\begin{defn}[!!{formula}s वरील विगमन]
प्रत्येक अण्वीय !!{formula} ला गुणधर्म~$P$ असेल आणि
!!{formula}-रचना-फलनांच्या खाली $P$ जपला जात असेल, तर प्रत्येक
!!{formula} ला गुणधर्म~$P$ असतो.
\end{defn}

\begin{explain}
या व्याख्येतून !!{formula}s वरील विगमनात्मक सिद्धतेची
“कृती” सुचते. प्रत्येक !!{formula} ला गुणधर्म~$P$
आहे हे दाखवायचे असल्यास पुढील पायऱ्या पाळा:\\

\textbf{विगमनाचा आधार.} $!A$ हे अण्वीय
!!{formula} असो. [\ldots] म्हणून $!A$ ला गुणधर्म~$P$ असतो.

\textbf{विगमनाची पायरी.} $!A$ आणि $!B$ ही दोन्ही
गुणधर्म~$P$ असलेली !!{formula}s असोत.

प्रकरण $\lnot$: [\ldots] म्हणून $\lnot !A$ ला गुणधर्म~$P$ असतो.

प्रकरण $\land$: [\ldots] म्हणून $!A \land !B$ ला गुणधर्म~$P$ असतो.

प्रकरण $\lor$: [\ldots] म्हणून $!A \lor !B$ ला गुणधर्म~$P$ असतो.

प्रकरण $\lif$: [\ldots] म्हणून $!A \lif !B$ ला गुणधर्म~$P$ असतो.

प्रकरण $\liff$: [\ldots] म्हणून $!A \liff !B$ ला गुणधर्म~$P$ असतो.

प्रकरण $\lforall[]$: [\ldots] म्हणून $\lforall[x] !A$ ला गुणधर्म~$P$ असतो.

प्रकरण $\lexists[]$: [\ldots] म्हणून $\lexists[x] !A$ ला गुणधर्म~$P$ असतो. \\

म्हणून प्रत्येक !!{formula} ला गुणधर्म~$P$ असतो.
\end{explain}

\end{document}
